% Figure: schematic of a heat pipe. Heat enters the evaporator end (left),
% boils the working fluid, the vapour streams down the hollow core to the
% condenser end (right) where it gives up its latent heat and condenses, and
% the wick returns the liquid by capillary action. No moving parts, no pump.
\begin{figure}[htbp]
\centering
\begin{tikzpicture}[font=\small]
  % outer sealed tube
  \draw[engblack, very thick] (0,0) rectangle (11,2.6);
  % wick layer (top and bottom inner walls)
  \fill[engblue!18] (0.08,0.08) rectangle (10.92,0.55);
  \fill[engblue!18] (0.08,2.05) rectangle (10.92,2.52);
  \node[engblue, font=\scriptsize, anchor=west] at (0.15,0.31) {wick (liquid return)};
  % vapour core
  \node[engred, font=\scriptsize] at (5.5,1.5) {vapour core};
  % vapour flow arrow (evaporator -> condenser)
  \draw[engred, line width=1.1pt, -{Stealth[length=3mm]}]
    (1.6,1.75) -- (9.4,1.75);
  % liquid return arrow (condenser -> evaporator) in the wick
  \draw[engblue, line width=1.1pt, -{Stealth[length=3mm]}]
    (9.4,0.31) -- (1.6,0.31);
  % section dividers
  \draw[enggray, dashed] (3.4,-0.5) -- (3.4,3.0);
  \draw[enggray, dashed] (7.6,-0.5) -- (7.6,3.0);
  % section labels
  \node[engorange, font=\scriptsize, anchor=south, align=center] at (1.7,2.7)
    {evaporator\\(heat in)};
  \node[enggray, font=\scriptsize, anchor=south, align=center] at (5.5,2.7)
    {adiabatic\\transport};
  \node[engblue, font=\scriptsize, anchor=south, align=center] at (9.3,2.7)
    {condenser\\(heat out)};
  % heat-in arrows at evaporator
  \foreach \x in {0.6,1.4,2.2,3.0} {
    \draw[engorange, line width=1pt, -{Stealth[length=2mm]}]
      (\x,-0.5) -- (\x,-0.02);
  }
  \node[engorange, font=\scriptsize, anchor=north] at (1.8,-0.6) {$\dot Q_{\text{in}}$};
  % heat-out arrows at condenser
  \foreach \x in {8.0,8.8,9.6,10.4} {
    \draw[engblue, line width=1pt, -{Stealth[length=2mm]}]
      (\x,2.62) -- (\x,3.1);
  }
  \node[engblue, font=\scriptsize, anchor=south] at (9.2,3.15) {$\dot Q_{\text{out}}$};
\end{tikzpicture}
\caption[Schematic of a heat pipe]{A heat pipe: a sealed tube lined with a
capillary wick and holding a working fluid at its saturation condition. Heat
entering the evaporator end boils the liquid; the vapour, driven by the tiny
pressure difference the boiling sets up, streams down the hollow core to the
condenser end, where it releases its latent heat and condenses; the wick pulls
the condensate back to the evaporator by capillary action. The device moves heat
along its length almost isothermally and with an effective conductance far above
any solid metal, because the transport rides the latent heat of the phase change
rather than a temperature gradient in a bar. It has no moving parts and needs no
external power, the reason heat pipes cool laptop processors, spacecraft
electronics, and the permafrost foundations of the trans-Alaska pipeline.}
\label{fig:vol04ch01:heat-pipe}
\end{figure}
